% UV Builder logo -- standalone TikZ building blocks.
% Compile with: pdflatex UVBuilderLogo.tex  (one page per tikzpicture)
\documentclass[tikz,border=24pt]{standalone}
\usepackage[T1]{fontenc}
\usepackage{tgpagella} % Palatino (TeX Gyre Pagella), scalable, for the logo text
\usepackage{mathrsfs}  % \mathscr: curly script L for the Lagrangians
\usetikzlibrary{arrows.meta,calc,fadings,decorations.pathmorphing,decorations.markings,decorations.text}

\definecolor{photon}{RGB}{230,120,20}
\definecolor{uvpurple}{RGB}{118,40,200}

% ---------------------------------------------------------------------------
% Template parameters (set them as options of a \pic, e.g. \pic[box size=3cm] {box};)
% ---------------------------------------------------------------------------
\tikzset{
  box size/.store in=\uvbBoxSize,           box size=4cm,
  box line width/.store in=\uvbBoxLineWidth, box line width=2pt,
  diagram width/.store in=\uvbDiagramWidth, diagram width=4cm,
  photon color/.store in=\uvbPhotonColor,   photon color=photon,
  % photon line width; default: same as the other lines of the diagram
  photon line width/.store in=\uvbPhotonWidth, photon line width=\pgflinewidth,
  hook size/.store in=\uvbHookSize,         hook size=1cm,
  % fermion line with an arrow along (fermion) or against (antifermion) the path
  fermion/.style={postaction={decorate},
    decoration={markings, mark=at position 0.55 with {\arrow{Stealth[scale=1.4]}}}},
  antifermion/.style={postaction={decorate},
    decoration={markings, mark=at position 0.55 with {\arrowreversed{Stealth[scale=1.4]}}}},
}

% ---------------------------------------------------------------------------
% box: rectangle of width `box width` x `box size` and height `box height` x
% `box size` (both fractions: 0 = flat, 1 = full size), drawn with
% `box line width`, centred on the pic origin, with both diagonals
% ---------------------------------------------------------------------------
\tikzset{
  box width/.store in=\uvbBoxWidth,   box width=1,
  box height/.store in=\uvbBoxHeight, box height=1,
  pics/box/.style={code={
    \pgfmathsetlengthmacro\uvbW{0.5*\uvbBoxWidth*\uvbBoxSize}
    \pgfmathsetlengthmacro\uvbH{0.5*\uvbBoxHeight*\uvbBoxSize}
    \draw[line width=\uvbBoxLineWidth] (-\uvbW,-\uvbH) rectangle (\uvbW,\uvbH);
    \draw[line width=\uvbBoxLineWidth] (-\uvbW,-\uvbH) -- (\uvbW,\uvbH);
    \draw[line width=\uvbBoxLineWidth] (-\uvbW,\uvbH) -- (\uvbW,-\uvbH);
  }},
}

% ---------------------------------------------------------------------------
% box stack: `box count` boxes stacked vertically, `box gap` apart (0 = sharing
% edges); the whole stack is centred on the pic origin
% ---------------------------------------------------------------------------
\tikzset{
  box count/.store in=\uvbBoxCount, box count=3,
  box gap/.store in=\uvbBoxGap,     box gap=0pt,
  pics/box stack/.style={code={
    \foreach \i in {1,...,\uvbBoxCount}
      \pic at (0,{(\i-(\uvbBoxCount+1)/2)*(\uvbBoxHeight*\uvbBoxSize+\uvbBoxGap)}) {box};
  }},
  % box row: boxes laid out horizontally, edge to edge (plus `box gap`), the
  % whole row centred on the pic origin. Box number `box joint` (0 = none) gets
  % width `box joint width` instead of `box width`.
  box joint/.store in=\uvbBoxJoint,            box joint=0,
  box joint width/.store in=\uvbBoxJointWidth, box joint width=1,
  pics/box row/.style={code={
    % total row width, in units of box size (gaps excluded)
    \pgfmathsetmacro\uvbRowT{\uvbBoxCount*\uvbBoxWidth
      + (\uvbBoxJoint>0 ? \uvbBoxJointWidth-\uvbBoxWidth : 0)}
    \foreach \i in {1,...,\uvbBoxCount}{
      \pgfmathsetmacro\uvbRowW{\i==\uvbBoxJoint ? \uvbBoxJointWidth : \uvbBoxWidth}
      % centre of box i: widths of the boxes before it + half its own width
      \pgfmathsetlengthmacro\uvbRowX{((\i-1)*\uvbBoxWidth
          + ((\uvbBoxJoint>0 && \i>\uvbBoxJoint) ? \uvbBoxJointWidth-\uvbBoxWidth : 0)
          + \uvbRowW/2 - \uvbRowT/2)*\uvbBoxSize
        + (\i-1-(\uvbBoxCount-1)/2)*\uvbBoxGap}
      \pic[box width=\uvbRowW] at (\uvbRowX,0) {box};
    }
  }},
}

% ---------------------------------------------------------------------------
% diagram: s-channel e+ e- -> gamma* -> e+ e-, total width `diagram width`,
% photon drawn in `photon color`, no labels; centred on the pic origin
% ---------------------------------------------------------------------------
\tikzset{
  pics/diagram/.style={code={
    \coordinate (-in-e)  at (-0.5*\uvbDiagramWidth,-0.3*\uvbDiagramWidth);
    \coordinate (-in-p)  at (-0.5*\uvbDiagramWidth, 0.3*\uvbDiagramWidth);
    \coordinate (-out-e) at ( 0.5*\uvbDiagramWidth, 0.3*\uvbDiagramWidth);
    \coordinate (-out-p) at ( 0.5*\uvbDiagramWidth,-0.3*\uvbDiagramWidth);
    \coordinate (-v1)    at (-\uvbDiagramWidth/6,0);
    \coordinate (-v2)    at ( \uvbDiagramWidth/6,0);
    \draw[fermion]     (-in-e) -- (-v1);
    \draw[antifermion] (-in-p) -- (-v1);
    \draw[fermion]     (-v2) -- (-out-e);
    \draw[antifermion] (-v2) -- (-out-p);
    \draw[color=\uvbPhotonColor, line width=\uvbPhotonWidth, decorate,
          decoration={snake, amplitude=0.025*\uvbDiagramWidth,
                      segment length=\uvbDiagramWidth/18}]
      (-v1) -- (-v2);
  }},
}

% ---------------------------------------------------------------------------
% hook: filled crane hook (ring on top, shank bending into a round bowl with a
% pointed tip). Size set by `hook size` (total height ~ 1.5 x hook size).
% Origin = top of the ring, so it can hang directly from a block or cable.
% Defines the coordinate (<name>-throat): bottom of the bowl, where loads hang.
% ---------------------------------------------------------------------------
\tikzset{
  pics/hook/.style={code={
    \begin{scope}[x=\uvbHookSize,y=\uvbHookSize]
      % ring
      \fill[even odd rule] (0,-0.17) circle[radius=0.17] (0,-0.17) circle[radius=0.08];
      % shank + bowl + tip, as a single filled outline
      \fill
        (-0.1,-0.3) -- (-0.1,-0.45)
        to[out=-90, in=90] (0.3,-0.9)                 % shank bends into the bowl
        arc[start angle=0, end angle=-210, radius=0.3] % inner edge of the bowl
        -- (-0.42,-0.5)                               % up to the tip
        -- (-0.47,-0.83)                              % back down the outside
        arc[start angle=160, end angle=360, radius=0.5]
        to[out=90, in=-90] (0.1,-0.45)
        -- (0.1,-0.3) -- cycle;
      \coordinate (-throat) at (0,-1.2);
    \end{scope}
  }},
}

% ---------------------------------------------------------------------------
% peak: lattice triangle (the crane's top mast), base width `box width` x
% `box size`, height `peak height` x `box size`, split in `peak levels`
% sections with X bracing; drawn with `box line width`. Origin = base centre.
% Defines the coordinate (<name>-apex).
% ---------------------------------------------------------------------------
\tikzset{
  peak height/.store in=\uvbPeakHeight, peak height=1.5,
  peak levels/.store in=\uvbPeakLevels, peak levels=3,
  pics/peak/.style={code={
    \pgfmathsetlengthmacro\uvbPW{0.5*\uvbBoxWidth*\uvbBoxSize}
    \pgfmathsetlengthmacro\uvbPH{\uvbPeakHeight*\uvbBoxSize}
    \begin{scope}[line width=\uvbBoxLineWidth, line join=round, line cap=round]
      \draw (-\uvbPW,0) -- (0,\uvbPH) -- (\uvbPW,0);
      \pgfmathtruncatemacro\uvbPN{\uvbPeakLevels-1}
      \foreach \k in {1,...,\uvbPN}{
        % section from level k-1 to level k: rung on top + X bracing
        \pgfmathsetlengthmacro\uvbYa{(\k-1)/\uvbPeakLevels*\uvbPH}
        \pgfmathsetlengthmacro\uvbYb{\k/\uvbPeakLevels*\uvbPH}
        \pgfmathsetlengthmacro\uvbXa{(1-(\k-1)/\uvbPeakLevels)*\uvbPW}
        \pgfmathsetlengthmacro\uvbXb{(1-\k/\uvbPeakLevels)*\uvbPW}
        \draw (-\uvbXb,\uvbYb) -- (\uvbXb,\uvbYb);
        \draw (-\uvbXa,\uvbYa) -- (\uvbXb,\uvbYb);
        \draw (\uvbXa,\uvbYa) -- (-\uvbXb,\uvbYb);
      }
    \end{scope}
    \coordinate (-apex) at (0,\uvbPH);
  }},
}

\begin{document}

% === crane ===================================================================
% --- tower and jib --------------------------------------------------------------
\def\stackCount{9}    % number of boxes in the vertical stack (tower)
\def\stackWidth{0.9}  % width of the tower boxes: 0 = flat, 1 = square
\def\rowCount{12}     % number of boxes in the horizontal row (jib)
\def\rowStart{3}      % which jib box (1 = leftmost) sits on the tower
\def\rowHeight{0.7}   % height of the jib boxes: 0 = flat, 1 = square
\def\latticeWidth{7pt}
% --- top of the crane -------------------------------------------------------------
\def\peakHeight{1.6}  % height of the triangle on top, in box sizes
\def\peakLevels{3}    % number of braced sections in the triangle
\def\tipLength{1.2}   % length of the pointed nose at the end of the jib, in box sizes
\def\cableWidth{4pt}  % line width of all cables (pendants, hoist, slings)
% --- counterweight and cab (sizes in box sizes) ---------------------------------
\def\weightWidth{0.8}   % counterweight, hanging past the short end of the jib
\def\weightHeight{1.6}
\def\weightGap{0.2}     % gap between the jib and the counterweight (bridged)
\def\cabWidth{1.9}      % cab: trapezoid centred on the tower, under the jib
\def\cabBottom{0.75}    % bottom width as a fraction of the top width
\def\cabHeight{0.9}
% --- trolley, hook and load ----------------------------------------------------
\def\trolleyBox{8}      % jib box the trolley hangs under (must be > \rowStart)
\def\trolleyWidth{0.8}  % in box sizes
\def\trolleyHeight{0.16}
\def\cableLength{0.5}   % hoist cable length, trolley to hook block, in box sizes
\def\hookSize{0.6}      % in box sizes
\def\slingDrop{2.6}     % vertical drop of the slings from the hook to the diagram
% load = Feynman diagram (e+ e- -> gamma -> e+ e-); the slings hold the two
% vertices at the ends of the photon line
\def\loadWidth{8.4}     % diagram width, in box sizes
\def\loadLineWidth{13pt}
\def\loadPhotonWidth{13pt}  % the photon line only
\def\loadPhotonColor{uvpurple}
% curly Ls (Lagrangians) under the photon, between the two lower legs
\def\lagCount{3}
\def\lagSize{0.9}        % font size, in box sizes
\def\lagSpacing{1.3}     % distance between neighbouring Ls, in box sizes
\def\lagDrop{0.17}       % how far below the photon line, in diagram widths
\def\lagDip{0.3}         % extra drop of the middle L(s), in box sizes (they sit on an arc)
\def\lagColor{black}
% UV spotlight: cone from the trolley down past the diagram, purple fading out
\def\lightColor{uvpurple}
\def\lightOpacity{0.55}  % opacity at the top of the cone
\def\lightSpread{1.1}    % width at the bottom, as a fraction of the diagram width
\def\lightReach{0.35}    % how far below the diagram's centre the cone ends, in diagram widths
% --- frame circle and logo text --------------------------------------------------
\def\circleX{4.95}       % centre of the circle, in box sizes (tower centre = origin)
\def\circleY{0.4}
\def\circleR{6.6}        % radius, in box sizes
\def\circleWidth{30pt}
\def\circleGap{0.15}     % blank margin where the circle meets the crane, in box sizes
\def\towerGap{0.2}       % blank margin between the bottom of the tower and the circle
\def\logoText{|\uvbLogoFont\color{uvpurple}|UV |\uvbLogoFont\color{black}|Builder}
\def\logoTextSize{1.3}   % font size, in box sizes
\def\logoTextGap{0.3}    % space between the circle and the text (inside), in box sizes
\def\logoTextSpan{60}    % half of the angle (degrees) the text may spread over

\begin{tikzpicture}[line join=round]
  % ---- key positions (the tower is centred on the origin) --------------------
  \pgfmathsetlengthmacro\s{\uvbBoxSize}
  \pgfmathsetlengthmacro\halfW{\stackWidth/2*\s}
  \pgfmathsetlengthmacro\jibBot{\stackCount/2*\s}
  \pgfmathsetlengthmacro\jibTop{(\stackCount/2+\rowHeight)*\s}
  \pgfmathsetlengthmacro\jibLeft{-\halfW-(\rowStart-1)*\s}
  \pgfmathsetlengthmacro\jibRight{\halfW+(\rowCount-\rowStart)*\s}
  \pgfmathsetlengthmacro\trolleyX{\halfW+(\trolleyBox-\rowStart-0.5)*\s}
  \pgfmathsetlengthmacro\trolleyBot{\jibBot-\trolleyHeight*\s}
  \pgfmathsetlengthmacro\blockTop{\trolleyBot-\cableLength*\s}
  \pgfmathsetlengthmacro\loadW{\loadWidth*\s}
  \pgfmathsetlengthmacro\fdY{\blockTop-0.38*\s-1.2*\hookSize*\s-\slingDrop*\s} % diagram centre
  \pgfmathsetlengthmacro\weightRight{\jibLeft-\weightGap*\s}
  \pgfmathsetlengthmacro\weightLeft{\weightRight-\weightWidth*\s}
  \pgfmathsetlengthmacro\weightTop{\jibTop+0.5*\latticeWidth}
  \pgfmathsetlengthmacro\noseX{\jibRight+\tipLength*\s}
  \pgfmathsetlengthmacro\apexY{\jibTop+\peakHeight*\s}

  % ---- frame circle, with blanks where the crane crosses it ------------------------
  % The circle is drawn first; then white "halo" shapes (crane parts, padded by
  % \circleGap) erase it wherever it would touch the crane; the crane is drawn on
  % top afterwards. The tower is instead cut off at the circle (see below).
  \coordinate (cc) at ({\circleX*\s},{\circleY*\s});
  \pgfmathsetlengthmacro\g{\circleGap*\s}
  \draw[line width=\circleWidth] (cc) circle[radius=\circleR*\s];
  % the bounding box ignores line widths: make room for the outer half of the stroke
  \path (cc) circle[radius=\circleR*\s+0.5*\circleWidth];
  \begin{scope}[white]
    % one silhouette around counterweight, jib, nose, triangle on top and both
    % pendant cables, so the circle gets a single continuous gap there
    \fill ({\weightLeft-\g},{\weightTop-\weightHeight*\s-\g})
      -- ({\weightLeft-\g},{\weightTop+\g})
      -- ({\jibLeft-\g},{\jibTop+\g+\cableWidth})
      -- (0,{\apexY+0.12*\s+\g})
      -- ({\jibRight+\g},{\jibTop+\g+\cableWidth})
      -- ({\noseX+2*\g},\jibBot)
      -- ({\noseX},{\jibBot-\g})
      -- ({\weightRight+\g},{\jibBot-\g})
      -- ({\weightRight+\g},{\weightTop-\weightHeight*\s-\g})
      -- cycle;
    % upper half of the tower (the lower half is cut off inside the circle instead)
    \fill ({-\halfW-\g},{\circleY*\s}) rectangle ({\halfW+\g},\jibBot);
    % cab under the jib (it sticks out below the silhouette above)
    \filldraw[line width=2*\g]
      ({-0.5*\cabWidth*\s},\jibBot) -- ({-0.5*\cabBottom*\cabWidth*\s},{\jibBot-\cabHeight*\s})
      -- ({0.5*\cabBottom*\cabWidth*\s},{\jibBot-\cabHeight*\s}) -- ({0.5*\cabWidth*\s},\jibBot)
      -- cycle;
    % the four outer legs of the diagram
    \draw[line width=2*\g+\loadLineWidth, line cap=round]
      ({\trolleyX-0.5*\loadW},{\fdY+0.3*\loadW}) -- ({\trolleyX-\loadW/6},\fdY)
      -- ({\trolleyX-0.5*\loadW},{\fdY-0.3*\loadW})
      ({\trolleyX+0.5*\loadW},{\fdY+0.3*\loadW}) -- ({\trolleyX+\loadW/6},\fdY)
      -- ({\trolleyX+0.5*\loadW},{\fdY-0.3*\loadW});
  \end{scope}

  % ---- logo text, following the bottom of the circle --------------------------------
  % The text sits inside the circle; its baseline runs left to right along an arc,
  % so the letters stand upright with their tops towards the centre.
  \pgfmathsetlengthmacro\fs{\logoTextSize*\s}
  \def\uvbLogoFont{\rmfamily\bfseries\fontsize{\fs}{\fs}\selectfont}
  \pgfmathsetlengthmacro\textR{\circleR*\s-0.5*\circleWidth-\logoTextGap*\s}
  \path[decorate, decoration={text along path, text align=center,
                              text={\logoText}}]
    ($(cc)+({270-\logoTextSpan}:\textR)$)
    arc[start angle={270-\logoTextSpan}, end angle={270+\logoTextSpan}, radius=\textR];

  % ---- tower (cut off inside the circle, leaving the same gap) + jib ------------
  \begin{scope}
    % only the lower half of the tower is cut at the circle: clip to the circle
    % plus everything above the circle's centre (both paths counter-clockwise, so
    % the clip is their union)
    \clip (cc) circle[radius=\circleR*\s-0.5*\circleWidth-\towerGap*\s]
          ({-\halfW-\s},{\jibTop+\s}) rectangle ({\halfW+\s},{\circleY*\s});
    \pic[box count=\stackCount, box width=\stackWidth, box line width=\latticeWidth]
      {box stack};
  \end{scope}
  \pic[box count=\rowCount, box height=\rowHeight, box line width=\latticeWidth,
       box joint=\rowStart, box joint width=\stackWidth]
    at ({((\rowCount+1)/2-\rowStart)*\s},{(\stackCount+\rowHeight)/2*\s}) {box row};

  % ---- triangle on top -----------------------------------------------------------
  \pic (peak) [box width=\stackWidth, box line width=\latticeWidth,
               peak height=\peakHeight, peak levels=\peakLevels]
    at (0,\jibTop) {peak};

  % ---- pointed nose at the end of the jib ----------------------------------------
  \draw[line width=\latticeWidth]
    (\jibRight,\jibTop) -- ({\jibRight+\tipLength*\s},\jibBot) -- (\jibRight,\jibBot);

  % ---- pendant cables from the apex to both ends of the jib ----------------------
  \draw[line width=\cableWidth] (\jibLeft,\jibTop) -- (peak-apex) -- (\jibRight,\jibTop);
  \fill (peak-apex) circle[radius=0.12*\s];   % pulley at the apex

  % ---- counterweight past the short end of the jib, bridged by the two chords ----
  \draw[line width=\latticeWidth]
    (\weightRight,\jibTop) -- (\jibLeft,\jibTop)
    (\weightRight,\jibBot) -- (\jibLeft,\jibBot);
  \fill[rounded corners=0.05*\s]
    ({\weightRight-\weightWidth*\s},{\jibTop+0.5*\latticeWidth})
    rectangle (\weightRight,{\jibTop+0.5*\latticeWidth-\weightHeight*\s});

  % ---- cab under the jib, centred on the tower, rounded bottom corners ------------
  \pgfmathsetlengthmacro\cabTop{0.5*\cabWidth*\s}
  \pgfmathsetlengthmacro\cabBot{0.5*\cabBottom*\cabWidth*\s}
  \fill (-\cabTop,\jibBot)
        [rounded corners=0.15*\s]
        -- (-\cabBot,{\jibBot-\cabHeight*\s})
        -- (\cabBot,{\jibBot-\cabHeight*\s})
        -- (\cabTop,\jibBot)
        [sharp corners] -- cycle;

  % ---- UV spotlight (drawn before the trolley, hoist and diagram, which sit on top
  % of it); it starts halfway up the trolley so it comes out from behind it ---------
  \begin{scope}
  \clip (cc) circle[radius=\circleR*\s-0.5*\circleWidth];
  \fill[\lightColor, fill opacity=\lightOpacity, path fading=south]
    ({\trolleyX-0.45*\trolleyWidth*\s},{\jibBot-0.5*\trolleyHeight*\s})
    -- ({\trolleyX+0.45*\trolleyWidth*\s},{\jibBot-0.5*\trolleyHeight*\s})
    -- ({\trolleyX+0.5*\lightSpread*\loadW},{\fdY-\lightReach*\loadW})
    -- ({\trolleyX-0.5*\lightSpread*\loadW},{\fdY-\lightReach*\loadW})
    -- cycle;
  \end{scope}

  % ---- trolley (centred under its jib box, slightly rounded bottom) --------------
  \fill ({\trolleyX-0.5*\trolleyWidth*\s},\jibBot)
        [rounded corners=0.03*\s]
        -- ({\trolleyX-0.5*\trolleyWidth*\s},\trolleyBot)
        -- ({\trolleyX+0.5*\trolleyWidth*\s},\trolleyBot)
        -- ({\trolleyX+0.5*\trolleyWidth*\s},\jibBot)
        [sharp corners] -- cycle;

  % ---- hoist cables, hook block, hook ---------------------------------------------
  \foreach \dx in {-0.12,0.12}
    \draw[line width=\cableWidth] ({\trolleyX+\dx*\s},\trolleyBot)
                                  -- ({\trolleyX+\dx*\s},\blockTop);
  \fill[rounded corners=0.06*\s]
    ({\trolleyX-0.25*\s},\blockTop) -- ({\trolleyX+0.25*\s},\blockTop)
    -- ({\trolleyX+0.18*\s},{\blockTop-0.4*\s}) -- ({\trolleyX-0.18*\s},{\blockTop-0.4*\s})
    -- cycle;
  \pic (hook) [hook size=\hookSize*\s] at (\trolleyX,{\blockTop-0.38*\s}) {hook};

  % ---- Feynman diagram, held by the two photon vertices ---------------------------
  \pic (fd) [diagram width=\loadW, photon color=\loadPhotonColor,
             photon line width=\loadPhotonWidth,
             line width=\loadLineWidth, line cap=round]
    at ($(hook-throat)+(0,{-\slingDrop*\s})$) {diagram};
  \draw[line width=\cableWidth, line cap=round]
    (fd-v1) -- (hook-throat) -- (fd-v2);

  % ---- Lagrangians under the photon ------------------------------------------------
  % (math fonts don't follow \fontsize, so the 10pt symbol is scaled instead)
  \pgfmathsetmacro\lagScale{\lagSize*\s/10pt}
  \foreach \k in {1,...,\lagCount}
    \node[text=\lagColor, inner sep=0pt, scale=\lagScale]
      at ({\trolleyX+(\k-(\lagCount+1)/2)*\lagSpacing*\s},
          {\fdY-\lagDrop*\loadW
           -\lagDip*\s*(1-((\k-(\lagCount+1)/2)/max((\lagCount-1)/2,0.5))^2)})
      {$\mathscr{L}$};
\end{tikzpicture}

\end{document}
